% Français 9115, batch 27 — Gallica views 521–540.
% Pass: Opus 5.5 (claude-opus-5-5), 2026-09-22 — first pass, unchecked against the views by a human.
%
% What the twenty views hold. Every one of the twenty views carries text:
% ten leaves written on both sides, in her regular fair hand, corrected
% between the lines (struck words, words added above, a few passages
% struck line by line, one paragraph cancelled by a large pen cross). It is
% one continuous text, paginated by her (111) to (130) at the head, centre,
% in parentheses, none struck; the layout is that of the prize memoir on
% vibrating elastic surfaces read in batches 20–22 (pages (1)–(34) at
% v. 398–440): both sides written, the ink number at the top right on the
% odd-paginated side only. Whether pages (35)–(110), which fall in batches
% 23–26, join the two runs into one memoir is for those batches to say.
% The subject throughout is rectangular plates after Chladni's « traité
% d'acoustique ». No microfilm target, no repeated exposure, no printed
% matter, no blank view.
%
%   v. 521      (111) the end of a discussion of Chladni's figures on
%               rectangular plates, explained by combining two cases, one
%               term being periodic: fig. 158 (« fig. 158 a »), a
%               « véritable ligne invariable de repos » through the middle;
%               she declines to consider his figures for plates whose side
%               ratio was « un peu alteré ». The first line is struck and
%               the sentence began before the batch — on the plate 5:4
%               and the cases 5|0 and 1|4, according to batch 26 (v. 520),
%               whose header also reads the formula (t^2m^2+s^2n^2)/t^2.
%   v. 522–523  (112)–(113) the plate whose sides are ::2:1: the cases 2/1
%               and 3/0 (figs. 173, « N.o 120 du traité d'acoustique »),
%               5/1 and 1/3 (figs. 174), two straight nodal lines crossing
%               at the centre; the theory « s'y applique parfaitement ». She
%               announces the comparison of the sounds of rectangular
%               plates, leaving out the « vibrations tournantes ». The text
%               stops three quarters down v. 523.
%   v. 524–526  (114)–(116) N.o 19, « comparaison des sons »: the numbers
%               to which the sounds are proportional given by a formula
%               read, with doubt, (t^2n^2+s^2m^2)/t^2; for the plate ::9:8 it
%               becomes (64m^2+81n^2)/64; a table of eighteen cases 2/2 to
%               6/3 (v. 525), then the same numbers divided by 8, the number
%               of the sound of Chladni's fig. 68 on the square plate
%               (v. 526).
%   v. 527–528  (117)–(118) Chladni's remark (« V. P. 161 ») that all the
%               rectangular plates sound about a semitone higher than the
%               square one; her conjecture that the glass of the square
%               plate differed (« quelque circonstance physique
%               inappréciable »), so she takes « sib3 + aigu » for fig. 68
%               instead of the « sib3 aigu » of his table P. 152; his table
%               P. 160 copied for the eighteen cases (notes and octaves);
%               the case 2/2 begun.
%   v. 529–532  (119)–(122) the eighteen comparisons, one paragraph each:
%               the interval heard, expressed as a ratio (15/8, 3, 9/2,
%               16/9, 8/3, 15/4, 16/3, 32/9, 32/5, 4, 6, 8, 50/9), against
%               the ratio of the theory, both brought to a common
%               denominator (960 and 985, 97 and 96, 2281 and 2304 …).
%   v. 533      (123) the conclusion: of the eighteen comparisons « il ne
%               se trouve qu'un seul cas » where theory and experiment
%               differ by more than a semitone (5/5, v. 532). Then the plate
%               ::2:1 again, « à cause de la simplicité du rapport »; a long
%               passage below is struck line by line.
%   v. 534–536  (124)–(126) the plate ::2:1: « En fesant donc s=2 t=1 »,
%               n^2+4m^2; the numbers and their ratios to 8 for eight cases;
%               Chladni's table P. 168; the eight comparisons (v. 535–536):
%               six within a semitone, two a little more than a tone off,
%               which did not prevent « M^r Chaldni » from treating the
%               figures as equivalent. A paragraph quoting Chladni's N.o 124
%               (sounds of the square plate « a peu près dans le rapport des
%               nombres 6, 15, 30 &c. », approaching « la série des nombres
%               naturels 1, 2, 3, 4 &c » as one diameter is reduced) is
%               cancelled by a large cross.
%   v. 537      (127) the figures of the periodic equations would also
%               appear on stretched surfaces (« surfaces tendues »), part of
%               them reversed for free sides; the sounds of the stretched
%               surface go as the simple numbers, those of the elastic
%               plate as their squares — « Euler a remarqué un semblable
%               rapport » between the lamina supported at both ends and the
%               string. She wished to experiment on plates with one or more
%               sides supported and could not.
%   v. 538–540  (128)–(130) a new section, titled and underlined: « De
%               l'effet de l'application d'un corps mou, aux côtés d'une
%               plaque rectangulaire élastique vibrante », preceded by a
%               mark read « § 6 » with doubt; N.o 20. A side clamped in a
%               grooved piece of wood lowered the sound instead of raising
%               it — the weight of the wood; she then loaded the sides with
%               « cire collante » and found the effect equivalent to a
%               fictitious lengthening of the plate. « experiences »: a
%               square glass plate of « un décimetre 6 millimetres », « 1
%               once 18 grains », the octave counted from the « la du
%               diapason »; its sound near « sol (1) »; with « 1 once 2 gros
%               de cire » on one side, more than a tone lower, « entre mi et
%               fa 1 »; the nodal line parallel to the loaded side at 45 mm
%               from it, 58 mm from the opposite side. View 540 ends on a
%               full stop, two thirds down the page.
%
% Seams. View 520 (outside the batch) ends « … 5 lignes nodales droites et
% parallelles entr'elles et » (the last « et » apparently struck); view
% 521 opens with a struck line « l'une de ces lignes passeroit par le
% milieu du plus » and goes on « et au plus grand côté dela plaque »: the
% sentence runs across 520/521, as batch 26 (written in parallel and read
% at the end of this pass) also says. At 540/541 nothing is cut: view 541 (page
% (131), ink 266, outside the batch) opens a new paragraph, « En laissant
% toujours un premier côté chargé de cire, j'en ai ajouté une egale
% quantité … », the same experiment continued.
%
% The numbers on the leaves, and why no \folio{} is written. (1) Her
% pagination, at the head, centre, in parentheses, not struck: (111) at
% v. 521 (the last figure blotted), then (112) to (130) at v. 522 to 540,
% one per view, without a gap. (2) The ink series at the top right, in the
% larger rounder hand: 256 (v. 521), 257 (523), 258 (525), 259 (527), 260
% (529), 261 (531), 262 (533), 263 (535), 264 (537), 265 (539) — one per
% leaf, always on the side her pagination makes odd, never on the even
% pages, which are the backs of the same leaves (the figure shows through,
% reversed, at v. 538); 266 follows at v. 541. Penned, as everywhere in
% this volume: no \folio{}. (3) Her paragraph numbers: N.o 19 (v. 524), N.o
% 20 (v. 538), and a « § 6 », read with doubt, before the title at v. 538.
% (4) References: Chladni's figures 68, 158, 158 a, 173, 174; his « traité
% d'acoustique » N.o 120 and N.o 124, pages 152, 160, 161, 168. Every
% number above was read on its view; none is inferred. None of these
% leaves can be Laubenbacher and Pengelley's « Manuscript C » (ff.
% 348r–349r): the subject is elasticity and acoustics, and the ink series
% here runs 256–265.
%
% Notation. Hers, kept. Chladni's cases, written as two figures separated
% by an oblique stroke (« 2/3 »: transverse and longitudinal nodal lines),
% are set $2|3$. Her « :: » of proportion is kept ($::9:8$). Notes are
% written as Chladni writes them, the name followed by the number of the
% octave, the sharp or flat set small above: set si$^{\flat}$3,
% fa$^{\sharp}$5, and « + aigu » as she writes it. Her points of
% multiplication (« 4.64 ») are kept. Words she underlines in the prose
% are set in \emph{}; letters underlined inside the mathematics keep
% \underline{}. The long s is transcribed as s; in the formula of v. 524
% and 533 it is the letter s, as v. 534 shows. Her tables are set as
% arrays, the dashes and dots between case and value as \ldots. Spelling
% is hers and is not flagged: « proportionel », « parallelles »,
% « equivallentes », « egals », « satisfesant », « fesant », « affoiblir »,
% « quarrée », « comma », « équivalans », « changemens », « instrumens »,
% « mouvemens », « differens », « rappellant », « raprocher », « Chaldni »
% once, and the run-together words « dela », « delaplaque », « lamême »,
% « lemême », « àdes », « adonné », « aucontraire », « àlathéorie »,
% « ladifference », « parconséquent », « desorte », « delà », « audessous »,
% « àprésent », « c'est àdire », « defort », « del'autre », « deses ».
% Where the leaf and the calculation disagree the leaf stands, with a note:
% a sign at v. 530 (« > que 16/3 » where the ratio is smaller) and the
% words « donné par l'experience » for « par la théorie » at v. 532.
%
% What could not be read, and what is offered with doubt. \ill{} stands
% twice: a short word struck with « ne pourrois apprécier » (v. 521); a
% short word struck with « de » before « plus d'un ton » (v. 540).
% \uncertain{} stands for « 18. » after « fig. 173 » (v. 522); the letters
% of the formula (t^2n^2+s^2m^2)/t^2 (v. 524, and the struck pair of v. 533);
% « l'on trouve » (struck) and « insensible » (v. 529); « Composé »
% (struck, v. 528); « depres » and « est » (v. 532); « ton » and « ton
% dela » (v. 533, struck); « Chaldni » (v. 536); « Je » (v. 537); « § 6 »
% (v. 538); the struck ending « oient » (v. 539). The case « 4/3 » of the
% last table of v. 534, blotted, is read from the first table of the same
% page and said so in a note.
%
% Nothing here was checked against any published edition of Chladni, of
% Euler or of Sophie Germain's papers.
\documentclass[11pt,a4paper]{article}
% The preamble shared by every transcription of the mathematicians' manuscripts in Gallica.
%
% It rests on one idea: **uncertainty compounds, it does not resolve.** A
% transcription that smooths over a doubtful word has destroyed the only thing
% it was made for — a reader must be able to tell, at every point, what was on
% the page and what was supplied. The apparatus macros below are therefore the
% critical apparatus, not decoration.
%
% The same commands are recognised by `scripts/render.mjs`, which produces the
% site's reading view, and by `scripts/tei.mjs`. A macro added here must be
% added there too, or rendering fails loudly — which is the wanted behaviour.
%
% Comments here are English, like the rest of the repository. The documents
% this preamble sets are French, and that is deliberate: see the skills.

\ProvidesPackage{germain}[2026/09/15 Transcriptions of mathematicians' manuscripts in Gallica]

\usepackage{amsmath,amssymb,amsthm}
\usepackage{mathrsfs}
% Kept from the parent preamble so the renderer's diagram support stays
% available; these manuscripts draw no commutative diagrams, and nothing here uses it.
\usepackage{tikz-cd}
\usepackage[english,french]{babel}
\usepackage{fontspec}

% --- Greek ------------------------------------------------------------------
% The text face stays fontspec's default, Latin Modern, which has no Greek:
% XeTeX drops every glyph it lacks with a line in the log and nothing on the
% page, and the Greek words the manuscripts quote vanished from the PDFs. So
% the Greek and Greek Extended blocks get a XeTeX character class of their own,
% and entering or leaving a run of it switches the family to CMU Serif — the
% same Computer Modern design, with polytonic Greek — and back. Only the
% family changes: a Greek word in \textit or \textbf keeps its shape and series.
% The transitions to and from every other class carry the tokens that class
% has towards an ordinary letter, so French spacing before « ; » or inside
% « » still holds after a Greek word. No grouping, so no brace or \endgroup
% can come out unbalanced; maths is left alone, since XeTeX inserts nothing
% there. Kept identical in `exercices.sty`.
\makeatletter
\newfontfamily\germainGreekfont{cmun}[
  Extension=.otf, NFSSFamily=germaingreek,
  UprightFont=*rm, ItalicFont=*ti, SlantedFont=*sl,
  BoldFont=*bx, BoldItalicFont=*bi, BoldSlantedFont=*bl]
\newXeTeXintercharclass\germain@greek
\def\germain@greekrange#1#2{%
  \@tempcnta=#1\relax
  \loop
    \XeTeXcharclass\@tempcnta=\germain@greek
    \ifnum\@tempcnta<#2\advance\@tempcnta\@ne\repeat}
\germain@greekrange{"0370}{"03FF}
\germain@greekrange{"1F00}{"1FFF}
\def\germain@greekprev{\rmdefault}
\def\germain@greekin{%
  \edef\germain@greekprev{\f@family}\fontfamily{germaingreek}\selectfont}
\def\germain@greekout{\fontfamily{\germain@greekprev}\selectfont}
\def\germain@greekjoin{%
  \XeTeXinterchartokenstate=\@ne
  \@tempcnta=\z@
  \loop
    \ifnum\@tempcnta=\germain@greek\else
      \XeTeXinterchartoks\germain@greek\@tempcnta=\expandafter{%
        \expandafter\germain@greekout\the\XeTeXinterchartoks\z@\@tempcnta}%
      \XeTeXinterchartoks\@tempcnta\germain@greek=\expandafter{%
        \the\XeTeXinterchartoks\@tempcnta\z@\germain@greekin}%
    \fi
    \ifnum\@tempcnta<4095\advance\@tempcnta\@ne\repeat}
% babel-french sets its own classes, and saves and restores the interchar
% state, each time a language is selected: join again after it does.
\addto\extrasfrench{\germain@greekjoin}
\addto\extrasenglish{\germain@greekjoin}
\AtBeginDocument{\germain@greekjoin}
\makeatother

\usepackage{geometry}
\geometry{a4paper,margin=25mm,marginparwidth=28mm}
\usepackage{marginnote}
\usepackage{xcolor}
\usepackage[normalem]{ulem}
\setlength{\emergencystretch}{3em}
\usepackage{hyperref}
\hypersetup{colorlinks=true,linkcolor=black,urlcolor=black,pdfborder={0 0 0}}

% --- Batch metadata ---------------------------------------------------------
% Taken as they stand by the rendering script, which shows them at the head of
% the reading view. They fix what the file claims to cover. The macro names are
% the parent project's, so that the scripts stay shared; \folder here means the
% volume's slug — `fr-9115` — and \pages counts Gallica views.

\newcommand{\folder}[1]{\gdef\grothFolder{#1}}
\newcommand{\batch}[1]{\gdef\grothBatch{#1}}
\newcommand{\pages}[2]{\gdef\grothFirst{#1}\gdef\grothLast{#2}}
\newcommand{\foldertitle}[1]{\gdef\grothFoldertitle{#1}}

% The holder's own shelfmark, printed as they write it: « Français 9115 ».
\newcommand{\shelfmark}[1]{\gdef\grothShelfmark{#1}}

% The dating, copied verbatim from the holder's catalogue — never inferred
% here. « XIXe siècle » is the BnF's; a pass that dates a leaf from its content
% says so in a \note{}, as its own inference.
\newcommand{\dating}[1]{\gdef\grothDating{#1}}

% The status notice, printed in the title block and repeated as a diagonal
% watermark on every page of the PDF. The manuscripts are in the public
% domain; what this line declares is not a right but a quality — an
% unverified first pass by a machine. The skills write the line; a file
% without it gets no watermark, so leaving it out is visible in the output.
\newcommand{\watermark}[1]{\gdef\grothWatermark{#1}}

\gdef\grothFolder{?}\gdef\grothBatch{}\gdef\grothFirst{?}\gdef\grothLast{?}%
\gdef\grothFoldertitle{}\gdef\grothShelfmark{}\gdef\grothDating{}\gdef\grothWatermark{}

% --- The résumé -------------------------------------------------------------
\newenvironment{resume}
  {\small\setlength{\parskip}{0.45em}}
  {\par\medskip\hrule\medskip}

% --- Keywords ---------------------------------------------------------------
% In English, closing the modernised reading's résumé: the modern vocabulary
% under which to search for what the leaves construct. The manifest extracts
% them to tag the volume — this line is the single source of the tags.
\newcommand{\keywords}[1]{\par\smallskip\noindent
  {\footnotesize\textsc{Keywords} --- \textit{#1}}\par}

% --- The critical apparatus -------------------------------------------------

% The Gallica view number, in the margin. This is the anchor that lets a
% disagreement over a formula be settled by looking at *one* image rather than
% a volume of 750. The site uses it to turn the facsimile as the transcription
% is scrolled. A view is one scanned image: usually one side of a leaf.
\newcommand{\page}[1]{%
  \marginnote{\footnotesize\color{gray!70}\textsf{v.\,#1}}%
  \label{page:#1}\ignorespaces}

% The BnF's pencilled foliation, where the leaf carries it and a pass can read
% it: « 198r ». This is what the literature cites — « ff. 198r–208v » — and
% the one line that turns a citation into a view. Recorded, never inferred:
% a folio a pass cannot read is not written.
\newcommand{\folio}[1]{%
  \marginnote{\footnotesize\color{gray!50}\textsf{f.\,#1}}\ignorespaces}

% The modernised reading's coarser counterpart to \page{}: which views a
% section is drawn from.
\newcommand{\pagerange}[2]{%
  \marginnote{\footnotesize\color{gray!70}\textsf{v.\,#1--#2}}%
  \ignorespaces}

% Illegible: never guessed.
\newcommand{\ill}{\textcolor{red!60!black}{[\,\ldots\,]}}

% A doubtful reading: the word is given, and flagged as uncertain.
\newcommand{\uncertain}[1]{\uline{#1}}

% An editorial addition — a missing letter, an implied word.
\newcommand{\add}[1]{\textcolor{blue!50!black}{[#1]}}

% Struck out by the writer. What was crossed out is part of the document.
\newcommand{\struck}[1]{\sout{#1}}

% The transcriber's note, on the state of the page or the mathematics.
\newcommand{\note}[1]{\par\smallskip{\small\color{gray!60!black}\textit{[#1]}}\par\smallskip}

% A marginal note of hers, distinct from the body of the text.
\newcommand{\marginal}[1]{\par\smallskip\noindent\hspace{1em}%
  {\small\color{gray!60!black}#1}\par\smallskip}

% --- Title ------------------------------------------------------------------
% The title block is French, like the documents themselves.

\AddToHook{shipout/background}{%
  \ifx\grothWatermark\empty\else
    \begin{tikzpicture}[remember picture, overlay]
      \node[rotate=52, scale=3.4, align=center, text=gray!18,
            font=\sffamily\bfseries]
        at (current page.center) {\grothWatermark};
    \end{tikzpicture}%
  \fi}

\AtBeginDocument{%
  \begin{center}
    {\large\bfseries Manuscrits de mathématiciens (Gallica) ---
      \ifx\grothShelfmark\empty\grothFolder\else\grothShelfmark\fi}\\[2pt]
    {\itshape\grothFoldertitle}\\[4pt]
    {\small vues \grothFirst--\grothLast
      \ifx\grothBatch\empty\else\ (lot \grothBatch)\fi}\\[2pt]
    \ifx\grothDating\empty\else
      {\small\color{gray!60!black}Datation du catalogue : \grothDating}\\[2pt]
    \fi
    \ifx\grothWatermark\empty\else
      {\small\bfseries\color{red!45!gray}\grothWatermark}\\[2pt]
    \fi
    {\footnotesize\color{gray!60!black}Transcription établie d'après les images IIIF de Gallica.
     Source gallica.bnf.fr / Bibliothèque nationale de France.\\
     Les lectures incertaines sont soulignées, les illisibles notées [\,\ldots\,],
     les ajouts entre crochets ; \textsf{v.} numérote les vues de Gallica,
     \textsf{f.} la foliotation au crayon de la BnF.}
  \end{center}
  \vspace{1em}}

\endinput


\folder{fr-9115}
\shelfmark{Français 9115}
\batch{27}
\pages{521}{540}
\dating{1801-1900}
\watermark{Lecture automatique\\non vérifiée}
\foldertitle{Papiers de Sophie GERMAIN. — Recueil de dissertations et problèmes mathématiques et physiques.}

\begin{document}

\note{Numérotation des feuillets. Chaque page porte en tête, au milieu,
entre parenthèses, sa pagination : (111) à la vue 521, puis (112) à (130)
aux vues 522 à 540, sans lacune et sans rature. Les feuillets sont écrits
des deux côtés, et seul le côté de pagination impaire porte en haut à
droite, à l'encre, d'une écriture plus ronde, un nombre de la série qui
court dans tout le volume : 256 (vue 521), 257 (523), 258 (525), 259
(527), 260 (529), 261 (531), 262 (533), 263 (535), 264 (537), 265 (539).
Cette série occupe la place de la foliotation de la Bibliothèque, mais
elle est tracée à la plume et non au crayon, et aucun « f. » n'est donc
porté. Tous ces nombres ont été lus sur la vue ; aucun n'est déduit.
Aucune vue de ce lot n'est blanche ni répétée.}

\page{521}

\struck{l'une de ces lignes passeroit par le milieu du plus}
et au plus grand côté dela plaque, \add{desorte} qu'une de ces lignes
passeroit par le milieu de \struck{cette} \add{la même} plaque et demême le
terme périodique, \struck{que seroit cos} \add{s'il étoit seul,} donneroit une figure
composée de quatre lignes droites parallelles au plus
petit côté dela \struck{même} plaque, et d'une seule dans
l'autre sens. \struck{car et} cette \add{dernière} ligne passeroit \add{donc} nécessairement par
le milieu dela même plaque, et seroit parconséquent
une véritable ligne invariable de repos. Les figures
158 montrent en effet l'existence de cette ligne qui
dans la fig. 158 \emph{a} reste seule droite, tandis que les
quatre autres sont sinueuses par l'influence du
terme périodique.

M\textsuperscript{r} Chladni adonné aussi plusieurs figures qui
repondent àdes cas où le rapport entre les cotes
des plaques ont été un peu alterés. Ainsi
je ne m'arreterai pas à les considérer, parce
\struck{e} que ne sachant pas dans quels cas ont été
operés les altérations des dimensions, je ne \struck{pourrois
apprécier} \ill{} ne puis apprécier leur influence.
\note{En tête, au milieu, entre parenthèses, « (111) », non barré, le dernier chiffre empâté ; en haut
à droite, à l'encre, d'une écriture plus ronde, « 256 ». La première
ligne, « l'une de ces lignes passeroit par le milieu du plus », est barrée
d'un trait continu ; la phrase reprend à la ligne suivante et commence
avant ce lot. « desorte » est écrit au-dessus de « qu'une » ; « la même »
au-dessus de « cette », barré. Un long trait horizontal passe sur « et
demême », comme ses traits de liaison, et n'est pas donné pour une
rature. À la ligne suivante, « que seroit cos » paraît biffé et « s'il
étoit seul, » est écrit au-dessus ; le « p » de « périodique » est chargé
d'un trait oblique. « car » et « et » sont barrés devant « cette » ;
« dernière » et « donc » sont écrits au-dessus de la ligne. Le « a »
souligné après « fig. 158 » désigne une des figures de Chladni. À
l'avant-dernière ligne, « ne pourrois / apprécier » est barré ; le mot
court qui suit, tracé comme « l'in », est touché par le même trait et ne
se lit pas.}

\page{522}

Les cas $2|1$ et $3|0$ donnent lieu, sur la plaque dont les
cotés sont entre eux $::2:1$ \struck{a des} aux fig. 173 \uncertain{18.} N.o 120 du
traité d'acoustique; on pourra voir, comme dans le
cas des fig. 158, que les nombres \struck{q} auxquels les sons
des deux termes separés doivent etre proportionnels,
different entr'eux de moins de 2 tons et $\frac{1}{2}$ ton,
limite dont cette difference approche, parceque
la figure qui appartiendroit à $2|1$ ne presenteroit
qu'un petit nombre de lignes nodales. On
voit aisément que dans le cas present il
ne doit pas y avoir de lignes invariables
de repos. Aussi les figures 173 n'ont elles
aucunes lignes nodales communes entr'elles.

A l'égard des fig. 174, \struck{On voit aisément}, que la
largeur doit etre, comme \struck{le dit} \add{l'observe} M\textsuperscript{r} Chladni, \add{(V. N.o 120)} un
peu plus grande que $\frac{1}{2}$ pour que les nombres
auxquels les sons qui \struck{accompa} repondent aux cas
$5|1$ et $1|3$ dont la combinaison explique ces
figures, soient egaux entr'eux. Je n'entreprendrai
pas de déterminer \struck{de combien doit etre} la mesure
de l'alteration, et je me contenterai de faire
\note{En tête, au milieu, « (112) », non barré ; aucun nombre en haut à
droite. C'est, semble-t-il, le dos du feuillet de la vue 521. Les cas
de Chladni sont écrits de deux chiffres séparés d'un trait oblique
(« 2/1 », « 3/0 », « 5/1 », « 1/3 ») ; ce trait est rendu par une barre
verticale. « $::2:1$ » est son signe de proportion. Après « fig. 173 »,
un groupe lu « 18. », peut-être un renvoi, précède « N.o 120 du traité
d'acoustique ». Après « les nombres », une lettre commencée est barrée.
« On voit aisément » est barré après « fig. 174, » ; « l'observe » est
écrit au-dessus de « le dit », barré ; « (V. N.o 120) » est écrit
au-dessus de la ligne, après « Chladni ». Le mot lu « largeur » commence
par une lettre tracée comme un « t ». Le « 5 » de « 5/1 » a la queue
descendante déjà relevée dans ce volume. Le « d » de « dont », avant « la
combinaison », est repassé ; « a » de « l'alteration » est surchargé. La
phrase se poursuit à la vue 523.}

\page{523}

observer que quelque soit le rapport des cotés, le cas de
$5|1$ et celui de $1|3$ donnent l'un et l'autre des figures
nodales dans lesquelles deux lignes de repos parallelles
à chacun des côtés delaplaque se croisent au centre
de \add{cette} plaque, \struck{desorte que les figures qui seront}
et que parconséquent les figures qui résultent dela
combinaison de ces deux cas doivent \struck{montrer} \add{presenter}, comme
elles \struck{montrent} \add{presentent} en effet, deux lignes nodales droites
\struck{qui sont de veritables} disposées comme je viens dele
dire, et qui sont de veritables lignes invariables de
repos.

Quoique je n'aie pas entrepris une explication détaillée
des differentes figures dont je viens deparler, j'espere
en avoir dit assez pour prouver que la théorie
\struck{est entièrement en état} s'y applique parfaitement.
Je vais àprésent \struck{entreprendre} \add{m'occuper} dela comparaison des
sons que rendent differentes plaques rectangulaires.
J'exclurai les cas qui appartiennent aux vibrations
tournantes, parceque, comme je l'ai souvent répété,
ces cas presentent des differences considerables dans
leur comparaison tant entre eux qu'avec \struck{les} ceux
\struck{ceux} où il y a \struck{bien} plusieurs périodes complettes
dans l'un et l'autre sens.
\note{En tête, au milieu, « (113) », non barré ; en haut à droite, à
l'encre, « 257 ». La phrase commencée à la vue 522 (« je me contenterai
de faire / observer ») se poursuit ici. « cette », devant « plaque », est
écrit plus haut que la ligne, dans un blanc laissé après « de ». Le
groupe « desorte que les figures qui seront » est barré d'un trait
continu, et la phrase reprend à la ligne suivante. « presenter » et
« presentent » sont écrits au-dessus de « montrer » et « montrent »,
barrés. « s'y » est écrit serré après le groupe barré « est entièrement
en état ». « m'occuper » est écrit au-dessus de « entreprendre », barré ;
« dela » n'est pas barré. « comparaison » porte à la fin un « s » barré
d'une petite croix : le pluriel est annulé et le singulier est donné.
Devant « ceux », « les » est barré ; au début de la ligne suivante un
second « ceux » est barré ; « bien » est barré devant « plusieurs ». Le
texte s'arrête aux trois quarts de la page, sur un point ; le bas du
feuillet est blanc.}

\page{524}

Dans ces comparaisons, je prendrai pour \struck{terme} premier terme
le son qui accompagne la fig 68 \struck{à laquelle} auquel
j'ai également comparé les differens sons qui accompagnent
les autres figures qui se montrent sur la plaque quarrée
je suivrai aussi, comme je l'ai fait alors, l'indication
de M\textsuperscript{r} Chladni \struck{suivant} d'après laquelle je regarderai
les figures comme \struck{formées} \add{equivallentes au système} delignes droites parallelles
entre elles, \add{et} \struck{a chacun des cotés} \struck{comme elles le seroient en effet
si elles appartenoient à l'intégrale (C).}
\add{parcequ'il ne s'agit ici que des cas compris dans les integrales periodiques}
et \struck{parconséquent comme appartenant à l'intégrale
périodique} \add{les \struck{sons} nombres auxquels les sons seront proportionnels}
\struck{(C), que donne la formule}
seront donnés par la formule \uncertain{$\frac{t^{2}n^{2}+s^{2}m^{2}}{t^{2}}$} \struck{pour
expression des nombres auxquels les sons des differens
cas doivent etre proportionnels.}

N.o 19. \emph{comparaison des sons}. Lorsque \add{les longueurs} des côtés dela plaque
sont entre elles $::9:8$ la formule qui donne les nombres
auxquels les sons sont proportionnels devient $\frac{64m^{2}+81n^{2}}{64}$
on trouvera donc en \struck{désignant} \add{notant}, avec M\textsuperscript{r} Chladni, par
$n/m$ la figure qui seroit composée de $\underline{n}$ lignes
transversales et de $m$ lignes longitudinales
\note{En tête, au milieu, « (114) », non barré ; aucun nombre en haut à
droite : c'est le dos du feuillet de la vue 523. La page est très
reprise. « terme » est barré devant « premier terme » ; « à laquelle »
devant « auquel ». Après « M\textsuperscript{r} Chladni », « suivant » est
barré. « equivallentes au système » est écrit au-dessus de « formées »,
barré. Au-dessus de « comme elles le seroient en effet », barré, un
ajout commence par « et », non barré, suivi de « a chacun des cotés »,
barré d'un trait ; la ligne suivante, « si elles appartenoient à
l'intégrale (C). », est barrée. « parcequ'il ne s'agit ici que des cas
compris dans les integrales periodiques » est écrit dans l'interligne, à
l'encre plus fine, au-dessus de la ligne qui commence par « et » ;
« parconséquent comme appartenant à l'intégrale » est barré, ainsi que
« périodique » au début de la ligne suivante ; « les nombres auxquels les
sons seront proportionnels » est écrit à sa suite, au-dessus de « (C),
que donne la formule », barré, avec « sons » barré devant « nombres ».
La formule n'est pas barrée ; ses lettres sont douteuses à cette place : la première lettre du numérateur et celle du dénominateur sont une même lettre barrée comme un $t$, la seconde lettre du numérateur est tracée comme un 1 ou un $l$ ; les exposants sont peu formés. Cette seconde lettre est lue $s$ long d'après la vue 534, où la même main écrit clairement « $s=2$ $t=1$ » et en tire $n^{2}+4m^{2}$ ; la lecture reste douteuse en bloc. Le reste de la phrase, de « pour » à « proportionnels. », est barré.
« comparaison des sons » est souligné ; « N.o 19 » est écrit en tête du
paragraphe. « les longueurs » est écrit au-dessus de la ligne, après
« Lorsque », et « des » est écrit sur un « les ». « notant » est écrit
au-dessus de « désignant », barré. Le $n$ de « $n$ lignes transversales »
est souligné. Le texte s'arrête aux deux tiers de la page, sans point ;
le bas du feuillet est blanc, et la phrase se poursuit à la vue 525.}

\page{525}

que le son de $2|2$ doit etre proportionel au nombre $\frac{4.64+4.81}{64}=\frac{145}{16}$
\[ \begin{array}{lcl}
2|3 & \ldots\ldots\ldots & \frac{4.64+9.81}{64}=\frac{985}{64} \\[4pt]
2|4 & \ldots\ldots\ldots & \frac{4.64+16.81}{64}=\frac{97}{4} \\[4pt]
2|5 & \ldots\ldots\ldots & \frac{4.64+25.81}{64}=\frac{2281}{64} \\[4pt]
3|2 & \ldots\ldots\ldots & \frac{9.64+4.81}{64}=\frac{225}{16} \\[4pt]
3|3 & \ldots\ldots\ldots & \frac{9.64+9.81}{64}=\frac{1305}{64} \\[4pt]
3|4 & \ldots\ldots\ldots & \frac{9.64+16.81}{64}=\frac{117}{4} \\[4pt]
3|5 & \ldots\ldots\ldots & \frac{9.64+25.81}{64}=\frac{2601}{64} \\[4pt]
4|2 & \ldots\ldots\ldots & \frac{16.64+4.81}{64}=\frac{337}{16} \\[4pt]
4|3 & \ldots\ldots\ldots & \frac{16.64+9.81}{64}=\frac{1753}{64} \\[4pt]
4|4 & \ldots\ldots\ldots & \frac{16.64+16.81}{64}=\frac{145}{4} \\[4pt]
4|5 & \ldots\ldots\ldots & \frac{16.64+25.81}{64}=\frac{3049}{64} \\[4pt]
5|2 & \ldots\ldots\ldots & \frac{25.64+4.81}{64}=\frac{481}{16} \\[4pt]
5|3 & \ldots\ldots\ldots & \frac{25.64+9.81}{64}=\frac{2329}{64} \\[4pt]
5|4 & \ldots\ldots\ldots & \frac{25.64+16.81}{64}=\frac{181}{4} \\[4pt]
5|5 & \ldots\ldots\ldots & \frac{25.64+25.81}{64}=\frac{3625}{64} \\[4pt]
6|2 & \ldots\ldots\ldots & \frac{36.64+4.81}{64}=\frac{657}{16} \\[4pt]
6|3 & \ldots\ldots\ldots & \frac{36.64+9.81}{64}=\frac{3033}{64}
\end{array} \]
ainsi, en se rappellant que le son qui accompagne la figure
68 doit, \struck{d'après la}, suivant la théorie, être proportionnel
au nombre 8, on trouvera, \struck{pour} que la difference
\struck{entre}
\note{En tête, au milieu, « (115) », non barré ; en haut à droite, à
l'encre, « 258 ». La première ligne achève la phrase de la vue 524
(« on trouvera donc […] par $n/m$ la figure […] / que le son de
$2|2$ doit etre proportionel au nombre »). Le tableau est écrit en
colonne : à gauche le cas, suivi d'une ligne de tirets, à droite la
formule ; les points qui séparent les facteurs sont les siens. Les
valeurs ont été relues sur la vue et non recalculées pour être lues ;
elles s'accordent avec la formule $\frac{64m^{2}+81n^{2}}{64}$ de la
vue 524. Au dénominateur de $\frac{337}{16}$ le 6 est écrit sur un autre
chiffre ; dans la dernière ligne, le 9 de « 9.81 » est repassé en
tache, et le 3 initial de « 36 » a la longue queue qui le fait prendre
pour un 9. « d'après la » est barré entre deux virgules. À la dernière
ligne, « pour » est barré ; le long trait qui court sous « que la
difference » est son trait de liaison et un soulignement, non une
rature ; « entre », seul à la ligne suivante, est barré, et la vue 526
reprend « la difference entre ». Le bas de la page est blanc.}

\page{526}

la difference entre les sons qui repondent
à la fig. 68 et à $2|2$ doit etre exprimée par $\frac{145}{128}$
\[ \begin{array}{ll}
\text{fig 68 et à } 2|3 \ldots\ldots\ldots & \frac{985}{512} \\[4pt]
\text{fig. 68 et à } 2|4 \ldots\ldots\ldots & \frac{97}{32} \\[4pt]
\text{fig. 68 et à } 2|5 \ldots\ldots\ldots & \frac{2281}{512} \\[4pt]
\text{fig. 68 et à } 3|2 \ldots\ldots\ldots & \frac{225}{128} \\[4pt]
\text{fig. 68 et à } 3|3 \ldots\ldots\ldots & \frac{1305}{512} \\[4pt]
\text{fig. 68 et à } 3|4 \ldots\ldots\ldots & \frac{117}{32} \\[4pt]
\text{fig. 68 et à } 3|5 \ldots\ldots\ldots & \frac{2601}{512} \\[4pt]
\text{fig. 68 et à } 4|2 \ldots\ldots\ldots & \frac{337}{128} \\[4pt]
\text{fig. 68 et } 4|3 \ldots\ldots\ldots & \frac{1753}{512} \\[4pt]
\text{fig. 68 et } 4|4 \ldots\ldots\ldots & \frac{145}{32} \\[4pt]
\text{fig. 68 et } 4|5 \ldots\ldots\ldots & \frac{3049}{512} \\[4pt]
\text{fig. 68 et } 5|2 \ldots\ldots\ldots & \frac{481}{128} \\[4pt]
\text{fig. 68 et } 5|3 \ldots\ldots\ldots & \frac{2329}{512} \\[4pt]
\text{fig. 68 et } 5|4 \ldots\ldots\ldots & \frac{181}{32} \\[4pt]
\text{fig. 68 et } 5|5 \ldots\ldots\ldots & \frac{3625}{512} \\[4pt]
\text{fig. 68 et } 6|2 \ldots\ldots\ldots & \frac{657}{128} \\[4pt]
\text{fig. 68 et } 6|3 \ldots\ldots\ldots & \frac{3033}{512}.
\end{array} \]
Avant de comparer les résultats que je viens de trouver
à ceux del'experience, il faut observer que, suivant la
\note{En tête, au milieu, « (116) », non barré ; aucun nombre en haut à
droite : c'est le dos du feuillet de la vue 525. La première ligne
reprend les derniers mots de la vue 525 (« que la difference / entre »,
où « entre » est barré). Chaque rapport est celui de la vue 525 divisé
par 8, le nombre du son de la fig. 68. Dans la ligne « fig. 68 et à
$4|2$ » les deux chiffres sont repassés à l'encre plus noire, le 4 sur
un autre chiffre ; dans la ligne suivante le 4 de « $4|3$ » l'est aussi.
À partir de cette ligne, « à » n'est plus écrit devant le cas. La phrase
finale se poursuit à la vue 527.}

\page{527}

remarque de M\textsuperscript{r} Chladni (V. P. 161), les sons qui répondent aux
mouvemens dela simple lame, sont plus aigus, sur la plaque dont
les longueurs des côtés sont entre elles $::9:8$, environ d'un demi
ton que sur la plaque quarrée. Cette différence \struck{est} à la verité,
\struck{peu considerable} \add{\struck{fort petite} n'est pas grande}; \struck{mais} et considerée dans un petit nombre
de cas, elle pourroit être négligée; mais sa constance la rend
remarquable: car toutes les plaques rectangulaires soumises à
l'experience présentent le même phénomène (V. les tables données
par M\textsuperscript{r} Chladni). Pour l'expliquer, il me semble naturel de
croire que quelque circonstance physique inappréciable, telle
qu'une différence de rigidité \struck{ou} du verre ou de son epaisseur
a donné pour la plaque quarrée des sons qui sont tous plus
graves d'environ un demi ton que si on avoit employé
le même verre dont on a taillé les plaques rectangulaires.

Cette conjecture est d'autant plus plausible qu'il suffit
que \struck{l'auteur physicien ait} se soit occupé à des époques differentes
de \struck{ces} ses experiences sur les plaques quarrées et sur les
plaques rectangulaires, pour qu'il \struck{n'}ait pû se procurer
du verre \struck{pris dans la même} exactement dela même
nature.

Conformément à l'observation que je viens de faire,
aulieu de prendre pour le son qui accompagne
la fig 68. si$^{\flat}$3 aigu que donne la table P. 152,
je prendrai au contraire si$^{\flat}$3 + aigu.
\note{En tête, au milieu, « (117) », non barré ; en haut à droite, à
l'encre, « 259 ». La phrase de la vue 526 (« il faut observer que,
suivant la / remarque de M\textsuperscript{r} Chladni ») se poursuit ici.
Dans « (V. P. 161) », le P est écrit sur une autre lettre, noircie. À la
troisième ligne, « les » est écrit sur un autre mot devant
« longueurs ». « est » est barré devant « à la verité » ; « peu
considerable » est barré, et au-dessus « fort petite », barré à son
tour, puis « n'est pas grande ». « mais » est barré devant « et
considerée ». Après « rigidité », un mot court est barré. « l'auteur
physicien ait » est barré devant « se soit occupé » ; « ces » corrigé en
« ses » ; « n' » barré devant « ait pû » ; « pris dans la même » barré.
Les notes sont écrites par leur nom suivi d'un chiffre, qui paraît être celui de l'octave ; le bémol est un petit signe en forme de « b » ou de « 6 » écrit au-dessus, entre le nom et le chiffre : « si$^{\flat}$3 ». La seconde fois le signe « + » précède « aigu ». Le texte s'arrête ici, aux quatre
cinquièmes de la page ; le bas du feuillet est blanc.}

\page{528}

\struck{\uncertain{Composé}} La table P. 160 du traité d'acoustique donne
pour le son de
\[ \begin{array}{lll}
2|2 & \ldots\ldots & \text{ut}\,4 \\
2|3 & \ldots\ldots & \text{la}\,4 + \text{aigu} \\
2|4 & \ldots\ldots & \text{fa}\,5 \\
2|5 & \ldots\ldots & \text{ut}.\,6 \\
3|2 & \ldots\ldots & \text{sol}^{\sharp}4 + \text{aigu} \\
3|3 & \ldots\ldots & \text{re}^{\sharp}5 + \text{aigu} \\
3|4 & \ldots\ldots & \text{la}\,5 + \text{aigu} \\
3|5 & \ldots\ldots & \text{re}^{\sharp}6 \\
4|2 & \ldots\ldots & \text{re}^{\sharp}5 + \text{aigu} \\
4|3 & \ldots\ldots & \text{sol}^{\sharp}\ldots\text{la}\,5 \\
4|4 & \ldots\ldots & \text{ut}^{\sharp}\ldots\text{re}\,6 \\
4|5 & \ldots\ldots & \text{fa}^{\sharp}6 \\
5|2 & \ldots\ldots & \text{si}^{\flat}5 \\
5|3 & \ldots\ldots & \text{ut}^{\sharp}\ldots\text{re}\,6 \\
5|4 & \ldots\ldots & \text{fa}.\,6 \\
5|5 & \ldots\ldots & \text{si}.^{\flat}6 \\
6|2 & \ldots\ldots & \text{mi}.\,6 \\
6|3 & \ldots\ldots & \text{fa}.^{\sharp}6
\end{array} \]
\struck{et} en attribuant \add{donc} si$^{\flat}$3 + aigu à la fig. 68, la difference entre ce ton
et celui de $2|2$ sera moindre qu'un ton majeur. \struck{En effet c'est a dire
que le rapport $\frac{145}{128}$ donné par la théorie est $<$ que $\frac{9}{8}$ d'une
quantité fort petite car on a $145 <$ que} Au contraire
le rapport donné par la théorie savoir $\frac{145}{128}$ est $>$ que $\frac{9}{8}$,
mais d'une quantité inappréciable puisque l'on a d'une part
145 et de l'autre 144.
\note{En tête, au milieu, « (118) », non barré ; aucun nombre en haut à
droite : c'est le dos du feuillet de la vue 527. Le premier mot de la
page, barré, est offert avec doute. Les notes sont écrites par leur nom
suivi d'un chiffre ; le dièse est un petit « $\sharp$ » et le bémol un petit
signe en forme de « b », écrits au-dessus entre le nom et le chiffre ;
le point qui suit parfois le nom (« ut. 6 », « fa. 6 », « mi. 6 ») est
gardé. « sol$^{\sharp}$ … la 5 », « ut$^{\sharp}$ … re 6 » : deux notes
réunies par des points, comme sur la feuille. Le « 5 » de « 5/2 », à la
treizième ligne du tableau, est écrit sur un autre chiffre. « et » est
barré au début de la phrase qui suit le tableau ; « donc » est écrit
au-dessus de « attribuant ». La phrase qui commence par « En effet »
est barrée jusqu'à « 145 $<$ que », avec « que », écrit au-dessus de
« le rapport » ; le signe « $<$ » est tracé comme un angle ouvert à
droite. Le bas de la page est blanc.}

\page{529}

La difference entre si$^{\flat}$3 + aigu et le son de $2|3$, est 5 tons et $\frac{1}{2}$ ton.
Elle est donc exprimée par $\frac{15}{8}$. Le rapport donné par la théorie
savoir $\frac{985}{512}$. differe très peu de celui là puisque l'on \add{trouve} en multipliant
par 512; d'une part 960 et de l'autre 985.

La difference entre si$^{\flat}$3 + aigu et le son de $2|4$, est un peu
moindre qu'un octave 3 tons et $\frac{1}{2}$ ton, c'est àdire $<$ que 3.
Aucontraire le rapport $\frac{97}{32}$ donné par la théorie est $>$ que 3,
mais d'une quantité inappréciable \add{à l'oreille} puisque l'on a d'une part
97 et de l'autre 96.

La difference entre si$^{\flat}$3 + aigu et le son de $2|5$, est
un peu moindre que 2 octaves et 1 ton, c'est àdire $<\frac{9}{2}$.
En effet le rapport $\frac{2281}{512}$ donné par la théorie, est $<$ que $\frac{9}{2}$
de moins d'un demi ton; car on \add{trouve} \struck{a} d'une part 2281 et de
l'autre 2304.

La difference entre si$^{\flat}$3 + aigu et le son de $3|2$, est environ
5 tons; elle est donc exprimée par $\frac{16}{9}$. En effet le rapport
$\frac{225}{128}$, \add{donné par la théorie} ne differe pas \struck{sensiblement} de $\frac{16}{9}$ \add{d'une quantité \uncertain{insensible} a l'oreille} puisque l'on \add{trouve} \struck{a} d'une
part 2025 et de l'autre 2048.

La différence entre si$^{\flat}$3 + aigu et le son de $3|3$, est environ
1 octave 2 tons et $\frac{1}{2}$ ton, c'est-à-dire qu'elle est exprimée par
$\frac{8}{3}$. Le rapport $\frac{1305}{512}$ \add{donné par la théorie} est $<$ que $\frac{8}{3}$; mais ladifference qui
\struck{est \uncertain{l'on trouve} etre $\frac{4096}{3915}$ est $<\frac{16}{15}$ car $\frac{3.1305}{8512}$ est $<\frac{25}{24}$
car $\frac{783}{512}$ est $<\frac{5}{3}$} entre ces deux rapports est bien moindre
qu'un demiton mineur; car \add{on a} $\frac{3.1305}{8.512}$ \struck{est} $<$ \add{que} $\frac{25}{24}$, $\frac{783}{512}<$ que $\frac{5}{3}$, $2349<$ que $2560$
\note{En tête, au milieu, « (119) », non barré ; en haut à droite, à
l'encre, « 260 ». Une page nouvelle, qui ne continue pas une phrase de
la vue 528 mais en poursuit le propos. Dans « 960 », le 6 est écrit sur
un autre chiffre ; dans « $2|4$ », le 4 aussi ; dans « $\frac{2281}{512}$ »
le dénominateur est repassé. « trouve » est écrit au-dessus de « l'on »
(première et troisième fois) et au-dessus de « a », barré (quatrième
fois) ; « à l'oreille » au-dessus de « inappréciable » ; « donné par la
théorie » au-dessus de la ligne, deux fois. Après « ne differe pas »,
« sensiblement » est barré, et « d'une quantité insensible a l'oreille »
est écrit au-dessus de « de $\frac{16}{9}$ puisque » ; le mot lu
« insensible » est offert avec doute, et « pas » n'est pas barré. Le
calcul biffé qui suit « ladifference qui » occupe une ligne et demie ;
son premier mot est « est », les deux suivants sont offerts avec doute ;
le dernier dénominateur de la ligne, écrit sur un autre nombre, est lu
24. Au premier des fractions finales, « on a » est écrit au-dessus de
« car », « est » est barré et « que » écrit au-dessus du signe $<$ ; le
3 initial des numérateurs « 3.1305 » a la longue queue qui le fait
prendre pour un 9. La dernière ligne atteint le bord du feuillet et
s'arrête sans point.}

\page{530}

La difference entre si$^{\flat}$3 + aigu et le son de $3|4$, est environ
1 octave 5 tons et $\frac{1}{2}$ ton; elle est donc exprimée par $\frac{15}{4}$.
Le rapport $\frac{117}{32}$ donné par la théorie diffère très peu de $\frac{15}{4}$; car
$\frac{120}{117}$ n'est guere plus grand que le \emph{comma}

La difference entre si$^{\flat}$3 + aigu et le son de $3|5$, est un peu
moindre que 2 octaves 2 ton et $\frac{1}{2}$ ton, c'est àdire $<$ que $\frac{16}{3}$.
Le rapport donné par la théorie, savoir $\frac{2601}{512}$ est en effet
$>$ que $\frac{16}{3}$, et la difference entre ces deux rapports \struck{est}
qui est exprimée par $\frac{8192}{7803}$, est environ $\frac{1}{2}$ ton.

La difference entre si$^{\flat}$3 + aigu et le son de $4|2$, est
environ 1 octave 2 tons et $\frac{1}{2}$ ton; elle peut donc être exprimée
par $\frac{8}{3}$. Le rapport $\frac{337}{128}$ donné par la théorie diffère en effet
très peu de $\frac{8}{3}$; car on trouve d'une part 1011 et de l'autre
1024.

La difference entre si$^{\flat}$3 + aigu et le son de $4|3$ \struck{est}, en prenant
un résultat moyen entre les deux notes qui appartiennent à ce
cas, \struck{de 1 octave et et 5 tons, environ}, environ, 1 octave et 5 tons,
c'est àdire qu'elle peut etre exprimée par $\frac{32}{9}$. Le rapport
$\frac{1753}{512}$ donné par la théorie diffère de $\frac{32}{9}$ d'un peu plus
d'un demiton; car on trouve d'une part 15777 et de
l'autre 16384. Cette difference qui est dans la limite de
\struck{l'erreur} la variation de l'experience seroit \struck{presque nulle} comme
on voit presque nulle si aulieu de prendre un résultat
moyen entre les deux sons on choisissoit celui qui s'accorde
le mieux avec l'experience.
\note{En tête, au milieu, « (120) », non barré ; aucun nombre en haut à
droite : c'est le dos du feuillet de la vue 529. « comma » est souligné.
Dans « 2 octaves 2 ton », le second 2 est écrit sur un autre chiffre.
Le signe devant « que $\frac{16}{3}$ », à la troisième ligne du deuxième
paragraphe, est tracé comme $>$ et est donné ainsi, alors que
$\frac{2601}{512}$ est plus petit que $\frac{16}{3}$, comme la phrase
précédente et le rapport $\frac{8192}{7803}$ le supposent. « est » est
barré à la fin de la ligne suivante. Dans « 1011 », les deux chiffres
du milieu sont empâtés. Au paragraphe suivant, « est » est barré après
« $4|3$ », et « de 1 octave et et 5 tons, environ » l'est aussi ;
« qu'elle » porte au-dessus de la fin un petit trait. « l'erreur » est
barré devant « la variation » ; « presque nulle » est barré à la fin de
la même ligne, « comme » ne l'est pas. Le bas de la page est blanc.}

\page{531}

La difference entre si$^{\flat}$3 + aigu et le son de $4|4$, est, en prenant
la note la plus basse qui est la plus favorable àlathéorie,
2 octaves 1 ton et près d'un demiton. Lerapport $\frac{145}{32}$ donné
par la théorie n'est que defort peu $>$ que $\frac{9}{2}$ qui exprime
un intervalle de 2 octaves et 1 ton; car on a d'une part
145 et de l'autre 144. ainsi \struck{en prenant} le rapport $\frac{145}{32}$ est
un peu $<$ que ne le veut l'experience; mais ladifference est
beaucoup moindre qu'un demiton, en prenant toute fois la
note la plus convenable.

La différence entre si$^{\flat}$3 + aigu et le son de $4|5$, est moins
de 2 octaves et 4 tons, c'est-àdire $<$ que $\frac{32}{5}$. En effet le rapport
$\frac{3049}{512}$, donné par la théorie est $<$ que $\frac{32}{5}$; car on trouve d'une
part 15245, et de l'autre 16384; desorte que les deux rapports
différent d'un peu plus d'un demi ton.

La différence entre si$^{\flat}$3 + aigu et le son de $5|2$, est un peu
moins de 2 octaves, parconséquent $<$ que 4. \struck{En effet} Le rapport
$\frac{481}{128}$ est effectivement $<$ que 4; car on trouve d'une part
481 et de l'autre 512. \struck{et} La différence entre les deux rapports
est donc environ un demi ton.

La différence entre si$^{\flat}$3 + aigu et le son de $5|3$, est, en prenant
la note la plus basse, 2 octaves 1 ton et près d'un demi ton.
Le rapport $\frac{2329}{512}$, donné par la théorie diffère àpeine de $\frac{9}{2}$
qui exprime un intervalle de 2 octaves et 1 ton; car on a d'une
part 2329 et de l'autre 2304; desorte que ladifference entre la
théorie et l'experience, est moindre qu'un demi ton en prenant à la verité la note la plus
convenable.
\note{En tête, au milieu, « (121) », non barré ; en haut à droite, à
l'encre, « 261 ». Une page nouvelle, qui poursuit la suite des cas
commencée à la vue 529. Ici le bémol de « si$^{\flat}$3 » est souvent
écrit sur la ligne, « si b 3 » ; il est rendu partout de même. Au
premier paragraphe, « en prenant » est barré après « ainsi ». Au
deuxième, « 4 tons » a le 4 tracé comme un « b », déjà relevé dans ce
volume ; le long trait qui court sous « En effet le rapport » est un
trait de liaison et n'est pas donné pour une rature. Au troisième, « En
effet » est barré devant « Le rapport », et « et » devant « La
différence ». La dernière ligne du feuillet, « théorie et l'experience,
[…] la note la plus », court jusqu'au bord droit, et le dernier mot,
« convenable. », est écrit au-dessous, à gauche, contre le bord inférieur
du feuillet.}

\page{532}

La difference entre si$^{\flat}$3 + aigu et le son de $5|4$, est un peu
moins de deux octaves 3 tons et $\frac{1}{2}$ ton, c'est àdire $<$ que 6.
En effet le rapport $\frac{181}{32}$ donné par la théorie est $<$ que 6,
et la difference deses deux nombres savoir $\frac{192}{181}$ est \struck{bien}
environ un demi ton.

La difference entre si$^{\flat}$3 + aigu et le son de $5|5$ est un peu
moindre que \struck{2} 3 octaves intervalle exprimé par 8.
Le rapport $\frac{3625}{512}$ est àlaverité $<$ que 8; mais l'excès d'un
de ces nombres sur l'autre étant \struck{d'environ 1 ton} \add{\uncertain{depres}} \add{d'un ton mineur} puisque
l'on \struck{trou} a d'une part 3625 et que l'on trouve de
l'autre 4096, il y a ici plus d'un demiton de difference
entre la théorie et l'experience.

La difference entre si$^{\flat}$3 + aigu et le son de $6|2$, est
un peu moindre que 2 octaves et 3 tons, c'est àdire $<$ que $\frac{50}{9}$
En effet le rapport $\frac{657}{128}$, donné par l'experience est
$<$ que $\frac{50}{9}$; car on trouve d'une part 5913 et de l'autre
6400. l'intervalle entre ces deux rapports \uncertain{est} environ un
demi ton.

La difference entre si$^{\flat}$3 + aigu et le son de $6|3$, est
un peu moindre que 2 octaves et 4 tons, c'est àdire $<$ que $\frac{32}{5}$
Le rapport $\frac{3033}{512}$, donné par la théorie est effectivement
$<$ que $\frac{32}{5}$; et ces deux rapports indiquent une distance
d'environ $\frac{1}{2}$ ton; car on trouve d'une part 15165 et 16384
del'autre.
\note{En tête, au milieu, « (122) », non barré ; aucun nombre en haut à
droite : c'est le dos du feuillet de la vue 531. Au premier paragraphe,
« bien » est barré à la fin de la quatrième ligne. Au deuxième, un
chiffre, lu 2, est barré et noirci devant « 3 octaves ». « d'environ 1
ton » est barré ; au-dessus sont écrits un mot lu « depres », offert
avec doute, puis « d'un ton mineur » ; « trou », mot commencé, est barré
devant « a d'une part ». Au troisième paragraphe, la leçon « donné par
l'experience » est bien celle de la feuille, là où les autres
paragraphes disent « donné par la théorie » ; le mot entre « rapports »
et « environ », tracé comme « en », est lu « est » avec doute. Au
dernier paragraphe, le 4 de « 4 tons » est écrit en gros sur un autre
chiffre. Le bas de la page est blanc.}

\page{533}

Je ne crois pas que l'on puisse esperer plus d'accord entre
une théorie et les faits auxquels elle s'applique, qu'il ne s'en
trouve ici entre \struck{les experiences} celle que je cherche à établir
et les experiences de M\textsuperscript{r} Chladni. On voit en effet que
parmi les 18 comparaisons que j'ai calculés entre les sons
dela plaque rectangulaire dont les côtés sont $::9:8$ et
le son corrigé de la fig. 68, il ne se trouve qu'un
seul cas où la difference entre les résultats dela théorie
et ceux de l'experience excède un demiton. \struck{Une
pareille précision justifie}.

Il me semble inutile d'effectuer, par rapport aux autres
plaques rectangulaires, des \struck{mêmes} comparaisons semblables à
celles que j'ai faites pour la première de ces plaques.
Cependant \struck{dans afin de faire voir que la théorie
s'applique également a tous les cas je vais calculer}, \add{à cause de la simplicité du rapport} \add{je m'occuperai encore de celle dont les côtés sont entr'eux $::2:1$} comparerai
\struck{dans tous ceux que M\textsuperscript{r} Chladni a soumis a l'experience
les résultats qu'il adonné a ceux que l'on peut tirer
des formules \uncertain{$\frac{t^{2}n^{2}+s^{2}m^{2}}{t^{2}}$, $\frac{t^{2}m^{2}+s^{2}n^{2}}{t^{2}}$} dont. La première s'applique
donne le nombre auquel le son doit être proportionnel
lorsqu'il y a $\underline{n}$ lignes transversales $\underline{m}$ lignes longitudinales
et \uncertain{ton} \add{\uncertain{ton dela}} seconde le nombre auquel le son doit etre proportion-
nel lorsqu'il y a $\underline{m}$ lignes transversales et $\underline{n}$ lignes
longitudinales.}
\note{En tête, au milieu, « (123) », non barré ; en haut à droite, à
l'encre, « 262 ». Une page nouvelle, qui conclut la comparaison des vues
529–532 et ouvre la suivante. « les experiences » est barré devant
« celle que » ; la phrase « Une pareille précision justifie. » est barrée
d'un trait ; « mêmes » est barré devant « comparaisons ». Le bas de la
page est un long passage barré ligne à ligne, de « dans afin de faire
voir » jusqu'à « longitudinales. » ; seuls échappent au trait
« Cependant », les deux ajouts de l'interligne et « comparerai », écrit
à la fin de la ligne « je vais calculer », barrée. Ces deux ajouts, d'une
écriture plus fine, sont écrits l'un au-dessus de « dans afin de
faire », l'autre au-dessus de « s'applique également » ; la phrase
gardée se lit donc : « Cependant, à cause de la simplicité du rapport je
m'occuperai encore de celle dont les côtés sont entr'eux ::2:1 »,
suivie de « comparerai », sans que la liaison soit écrite. Le « s » des
deux formules barrées est un s long, et la première lettre, barrée
comme un $t$, pourrait être un $f$ : lecture douteuse, comme celle de la
formule de la vue 524, qui a sans doute les mêmes lettres. À l'avant-dernière ligne barrée, le mot
qui suit « et », tracé comme « ton » ou « l'on », est offert avec doute ;
au-dessus est écrit, plus fin, un groupe lu « ton dela », douteux lui
aussi ; le sens attend « la seconde ». Un groupe de points d'encre sous la dernière ligne barrée n'est
pas de l'écriture. Le bas de la page est blanc.}

\page{534}

En fesant donc $s=2$ $t=1$, la formule qui donne lesnombres auxquels
les sons doivent etre proportionels est $n^{2}+4m^{2}$, desorte que le son
de $2|2$ doit etre proportionnel au nombre $4+4.4=20$
\[ \begin{array}{lcr}
\ldots\ 2|3 & \ldots\ldots\ldots & 4+4.9=40 \\
3|2 & \ldots\ldots\ldots & 9+4.4=25 \\
3|3 & \ldots\ldots\ldots & 9+4.9=45 \\
4|2 & \ldots\ldots\ldots & 16+4.4=32 \\
4|3 & \ldots\ldots\ldots & 16+4.9=52 \\
5|2 & \ldots\ldots\ldots & 25+4.4=41 \\
5|3 & \ldots\ldots\ldots & 25+4.9=61
\end{array} \]
En se rappellant donc que le son qui accompagne la fig 68 doit
etre proportionnel au nombre 8, on trouvera \struck{que} la difference
des sons qui repondent a cette fig. 68 et a $2|2$ doit etre exprimée par $\frac{5}{2}$
\[ \begin{array}{ll}
\text{fig. 68 et a } 2|3 \ldots\ldots\ldots & 5 \\[2pt]
\text{fig 68 et a } 3|2 \ldots\ldots\ldots & \frac{25}{8} \\[4pt]
\text{fig. 68 et a } 3|3 \ldots\ldots\ldots & \frac{45}{8} \\[4pt]
\text{fig. 68 et } 4|2 \ldots\ldots\ldots & 4 \\[2pt]
\text{fig. 68 et } 4|3 \ldots\ldots\ldots & \frac{13}{2} \\[4pt]
\text{fig. 68 et } 5|2 \ldots\ldots\ldots & \frac{41}{8} \\[4pt]
\text{fig 68 et } 5|3 \ldots\ldots\ldots & \frac{61}{8}
\end{array} \]
La table P. 168 du traité d'acoustique donne
pour le son de
\[ \begin{array}{lll}
2|2 & \ldots\ldots & \text{ut}\,5 \\
2|3 & \ldots\ldots & \text{ut}\,6 \\
3|2 & \ldots\ldots & \text{fa}^{\sharp}5 \\
3|3 & \ldots\ldots & \text{re}^{\sharp}6 \\
4|2 & \ldots\ldots & \text{si}\,5 \\
4|3 & \ldots\ldots & \text{fa}^{\sharp}6 \\
5|2 & \ldots\ldots & \text{re}^{\sharp}6 \\
5|3 & \ldots\ldots & \text{si}.^{\flat}6
\end{array} \]
\note{En tête, au milieu, « (124) », le 4 repassé ; aucun nombre en haut
à droite : c'est le dos du feuillet de la vue 533. La plaque dont les
côtés sont ::2:1, annoncée à la vue 533. Les lettres $s$ et $t$ sont
écrites ici clairement, le $s$ long : ce sont celles des formules barrées
de la vue 533 et, sans doute, de la vue 524 ; avec $s=2$, $t=1$ la
formule $\frac{t^{2}n^{2}+s^{2}m^{2}}{t^{2}}$ donne bien $n^{2}+4m^{2}$.
La première ligne du premier tableau commence par des points, sans cas
écrit devant ; le cas « $2|3$ » est écrit plus loin sur la ligne. Dans
la ligne « $4|2$ » du premier tableau, le dernier 4 de « 4.4 » est
repassé. « que » est barré, en tache, devant « la difference ». Après
« donne », à la ligne « La table P. 168 », un trait horizontal court
jusqu'à la marge. Dans le dernier tableau, le cas « $4|3$ » est écrit en
tache sur d'autres chiffres, un dièse jeté par-dessus ; il est lu d'après
le premier tableau, avec doute. Le bas de la page est blanc.}

\page{535}

On voit que l'intervalle entre $2|2$ et $2|3$, est exactement lemême
du côté dela théorie que du côté del'experience; ainsi il suffira
de comparer le son de $2|2$ a celui si$^{\flat}$3 + aigu. La différence
est 1 octave et 1 peu moins d'un ton. Le rapport $\frac{5}{2}$ est
donc trop grand d'un peu plus d'un ton que l'experience
ne le donne.

La difference entre si$^{\flat}$3 + aigu et le son de $3|2$, est
1 octave, et un peu moins de 4 tons c'est àdire $<$ que $\frac{16}{5}$.
Le rapport $\frac{25}{8}$ \add{donné par la théorie} est en effet $<$ que $\frac{16}{5}$ \struck{juste} \add{exactement} d'un \emph{comma}.

La difference entre si$^{\flat}$3 + aigu et le son de $3|3$, est
un peu moindre que 2 octaves 2 tons et $\frac{1}{2}$ ton, parconséquent
$<$ que $\frac{16}{3}$. Aucontraire $\frac{45}{8}$ est $>$ que $\frac{16}{3}$; mais \struck{d'une
quantité} la difference entre ces deux rapports est plus petite
qu'un demi ton majeur

La différence entre si$^{\flat}$3 + aigu et le son de $4|2$, est un
peu plus grande que 2 octaves; elle est donc à fort peuprès
exprimée par 4 qui est aussi le nombre indiqué par lathéorie

La difference entre si$^{\flat}$3 + aigu et le son de $4|3$, est un
peu moindre que 2 octaves et \struck{5} 4 tons, c'est àdire $<$ que $\frac{32}{5}$
Au contraire le rapport $\frac{13}{2}$ est $>$ que $\frac{32}{5}$, mais d'une quantité
inappréciable à l'oreille

La difference entre si$^{\flat}$3 + aigu et le son de $5|2$, est
un peu moindre que 2 octaves 2 tons et $\frac{1}{2}$ ton, c'est àdire
$<$ que $\frac{16}{3}$. En effet le rapport $\frac{41}{8}$ est un peu $<$ $\frac{16}{3}$
car on a d'une part 128 et de l'autre 123.
\note{En tête, au milieu, « (125) », non barré ; en haut à droite, à
l'encre, « 263 ». Une page nouvelle, qui poursuit, pour la plaque dont les
côtés sont ::2:1, la comparaison préparée à la vue 534. Au deuxième
paragraphe, une petite croix est écrite après « 1 octave » ; « donné par
la théorie » est écrit au-dessus de « est en effet » ; « exactement » est écrit au-dessus de « juste », barré ;
« comma » est souligné. Au troisième, « d'une quantité » est barré.
Au cinquième, un chiffre, lu 5, est barré devant « 4 tons ». Au dernier,
le signe devant $\frac{16}{3}$ est écrit sur un autre, en tache ; il est
lu $<$, comme le veulent le « En effet » et les nombres 128 et 123. Le
bas de la page est blanc.}

\page{536}

La différence entre si$^{\flat}$3 + aigu et le son de $5|3$, est un peu moindre
que 3 octaves, c'est àdire $<$ que 8. En effet le rapport $\frac{61}{8}$
\struck{est} donné par la théorie, est $<$ que 8 d'environ un demi
ton, car on a d'une part 61 et de l'autre 64.

On voit que sur les 8 comparaisons auxquels le cas présent donne
lieu, il y en a six dont l'accord avec la théorie est entièrement
satisfesant, puisque la difference entre le rapport qu'elle
indique et l'intervalle donné par l'experience est audessous
d'un demi ton. Mais on voit en même tems qu'il y a
deux cas où ladifference atteint son maximum puisqu'elle
est même un peu plus grande qu'un ton \add{mais} j'ai déja
remarqué qu'une pareille difference\add{, qui au reste est peu considerable,} n'a pas empêchée
M\textsuperscript{r} \uncertain{Chaldni} de regarder les figures auxquelles elles
appartenoient comme équivalentes. \struck{au reste je ne pretends
pas en deriver la cause et j'ai cru devoir}.

Après avoir donné les differentes tables des sons produits
sur les plaques rectangulaires M\textsuperscript{r} Chladni rassemble N.o 124
diverses reflexions générales sur la loi que suivent ces sons
on lit « que dans les vibrations où une ligne longitudinale est
« coupée par des transversales $1|1$, $2|1$, $3|1$ \&c. les sons d'une
« plaque quarrée sont a peu près dans le rapport des nombres
« 6, 15, 30 \&c. » Il ajoute que si l'un des diamètres étant
« peu a peu diminué, les sons se rapprochent de plus en
« plus de la série des nombres naturels 1, 2, 3, 4 \&c »
\note{En tête, au milieu, « (126) », le 6 repassé ; aucun nombre en haut
à droite : c'est le dos du feuillet de la vue 535. « est » est barré
devant « donné par la théorie ». « mais » est écrit au-dessus de la
ligne, entre « ton » et « j'ai » ; « , qui au reste est peu
considerable, » au-dessus de « difference n'a pas empêchée ». Le nom
est écrit ici, semble-t-il, « Chaldni », les lettres interverties ; il
est offert avec doute. Dans « équivalentes » une grosse tache couvre une
lettre. La fin du paragraphe, « au reste je ne pretends pas en deriver
la cause et j'ai cru devoir », est barrée d'un trait ; « deriver » est
douteux. Le dernier paragraphe, où elle cite le N.o 124 du traité de
Chladni, est cancellé tout entier par une grande croix à la plume, dite ici
plutôt que rendue par des ratures, comme aux vues 281 et 289 ;
le texte cité y est souligné, et les guillemets sont répétés en tête de
chaque ligne, comme elle le fait ailleurs ; « ces » est écrit sur un
autre mot devant « sons ». Le bas de la page est blanc.}

\page{537}

Parmi les figures qui se forment sur les differentes plaques rectangulaires,
\struck{il} on sent aisément que celles qui résultent des seules équations
périodiques, se montreroient egalement sur les surfaces tendues,
en observant pourtant que s'il s'agissoit d'une plaque rectangulaire
dont les côtés \struck{fussent} libres, \struck{les figures une} partie de ces
figures se trouveroit renversée sur la surface tendue, et qu'il
n'y auroit de ressemblance parfaite qu'entre les figures
dela même surface et \add{celles} dela plaque élastique appuyée
de tout côté. Il est egalement clair \struck{que}, d'après l'ordre
des équations differentielles qui conviennent à chacun de
ces corps sonores, que, tandis que les sons dela surface
tendue seront proportionnels à la suite des simples
nombres qui leur conviennent respectivement, les sons
dela plaque élastique seront proportionels aux
quarrés des mêmes nombres. Euler a \struck{fait} remarqué
un semblable rapport entre la simple lame
appuyée à ses deux extremités et la corde tendue.

J'aurois desiré faire quelques experiences sur les plaques élastiques
dont un ou plusieurs coté eussent été appuyés. \uncertain{Je}
n'ai pû atteindre le but que je me proposois; mais les
\add{tentatives que j'ai faites à cet egard m'ont donné}
\struck{j'ai eu} occasion d'observer et de mesurer un phénomène
\note{En tête, au milieu, « (127) », non barré ; en haut à droite, à
l'encre, « 264 ». Une page nouvelle : les plaques rectangulaires
comparées aux surfaces tendues. « il » est barré au début de la
deuxième ligne ; « fussent » est barré devant « libres » ; « les figures
une » est barré. « celles » est écrit au-dessus de « dela plaque ». « que »
est barré après « clair ». Dans « plaque élastique seront », le « s »
final de « plaques » est barré d'une petite croix : le singulier est
donné. Le nom d'Euler est écrit avec son E en forme de L. « fait » est
barré devant « remarqué ». Au second paragraphe, « un » est écrit sur un
autre mot ; à la fin de la ligne « appuyés. », une initiale noircie,
lue « Je », commence la phrase suivante. « tentatives que j'ai faites à
cet egard m'ont donné » est écrit dans l'interligne, à l'encre plus
fine, au-dessus de « j'ai eu », barré. La phrase se poursuit à la vue
538.}

\page{538}

à l'exposition duquel la classe me pardonnera peut-être
de m'arreter un instant

\uncertain{§ 6} \emph{De l'effet de l'application d'un corps mou, aux côtés d'une
plaque} \add{rectangulaire} \emph{élastique vibrante}.

N.o 20. Je cherchois a gêner le mouvement d'une des lignes de
limite d'une surface quarrée, en la serrant dans un morceau de
bois auquel j'avois pratiqué une rainure destinée à recevoir
un millimetre environ du côté de cette surface. J'esperois
produire \struck{je} ainsi pour cette ligne de limite, les effets
analytiques d'un simple appui. si j'eusse atteint ce but,
le son \struck{devoit} produit par une même figure devoit
devenir plus aigu que sur la plaque entièrement
libre. Bien loin delà, il arriva aucontraire que le
son étoit sensiblement plus grave que dans ce dernier
cas. Je ne tardai pas a penser que j'avois compliqué
l'effet que je voulois produire, d'un autre tout contraire
et que c'étoit au poids de ma pièce de bois que je
devois attribuer le singulier résultat de ma tentative.
Pour verifier ce soupçon, j'ai employé dela cire \struck{verte} collante
que j'ai distribuée uniformement, + dans toute l'étendue +
d'un ou de plusieurs côtés de differentes plaques rectangulaires
j'ai réduit ainsi l'effet de la gêne à affoiblir
\note{En tête, au milieu, « (128) », non barré ; aucun nombre en haut à
droite : c'est le dos du feuillet de la vue 537, dont le « 264 » se lit à
l'envers par transparence. Les deux premières lignes achèvent la phrase
de la vue 537 (« un phénomène / à l'exposition duquel »). Le titre est
souligné ; devant lui, un signe lu « § 6 », tracé comme « 16 », est
offert avec doute. « rectangulaire » est écrit au-dessus de
« élastique », qui n'est pas barré. « aux côtés » est écrit sur « au
côté ». « N.o 20 » suit le « N.o 19 » de la vue 524. « je », mot
commencé, est barré après « produire » ; « devoit » est barré devant
« produit ». L'initiale de « Bien » est repassée en tache. « verte »
paraît barré devant « collante ». « distribuée » porte à la fin un « s »
barré d'une petite croix : le singulier est donné. Les deux croix qui
encadrent « dans toute l'étendue » sont les siennes. La page s'arrête
sans point au bas du feuillet, et la phrase se poursuit à la vue 539.}

\page{539}

un peu le son, et a rendre parconséquent les figures nodales un
peu moins nettes qu'elles ne l'eussent été sur une surface non chargée
de cire. J'ai reconnu que les changemens produits par l'addition
du poids de cette même cire, étoient entièrement équivalans
à ceux que produiroient le prolongement effectif dela
plaque vibrante; c'est-àdire que le son \struck{est} \add{étoit} lemême que si
la plaque \struck{étoit} \add{eut été} réellement plus grande d'une portion
capable de faire équilibre àla quantité de cire ajoutée.
\struck{et que} Les lignes nodales changent\struck{\uncertain{oient}} \add{aussi} de place conformément
à ce prolongement fictif, de manière que la figure ne pouvoit
être supposée la même que sur la plaque non chargée,
qu'en ayant egard au prolongement fictif des lignes qui
la composent, sur la portion dont l'addition dela cire \struck{tenoit}
lieu.

J'ai fait un grand nombre d'experiences qui \struck{tendent} \add{donnent} toutes
\struck{a établir} le même résultat; mais je me contenterai d'en
rapporter un petit nombre, et je choisirai les plus aisées
à répeter, parcequ'elles suffisent pour mettre le phénomène
dans tout son jour.

\emph{experiences}

J'ai pris une plaque de verre quarrée dont le côté avoit
un décimetre 6 millimetres, cette plaque pesoit 1 once 18 grains;
et, en prenant pour premier octave celui auquel appartient
\note{En tête, au milieu, « (129) », non barré ; en haut à droite, à
l'encre, « 265 ». La première ligne achève la phrase de la vue 538 (« à
affoiblir / un peu le son »). « étoit » est écrit au-dessus de « est »,
barré ; « eut été » au-dessus de « étoit », barré. « et que » est barré
devant « Les lignes nodales » ; la fin de « changent » est surchargée,
et les lettres barrées sont offertes avec doute ; « aussi » est écrit
au-dessus. « tenoit » est barré devant « lieu » ; « tendent » devant
« toutes », avec « donnent » écrit au-dessus ; « a établir » est barré.
« experiences », souligné, est un intertitre au milieu de la ligne. Le
« 6 » de « 6 millimetres » est tracé comme un « b ». La page s'arrête sans
point au bas du feuillet, et la phrase se poursuit à la vue 540.}

\page{540}

le \emph{la} du diapason dont on se sert ordinairement pour
accorder les instrumens, le son dela figure composée de
deux lignes nodales parallelles, l'une a l'axe des $\underline{x}$ l'autre
a celui des $y$, s'est trouvé très près de sol (1) mais
un peu plus bas que cette note.

Après avoir collé à une des lignes d'extrémités dela
même plaque 1 once 2 gros de cire egalement distribuée
sur tous les points de cette ligne, j'ai continué de
tirer le son, comme avant l'addition dela cire, d'un
des angles non chargés dela plaque. Le son s'est
trouvé plus bas \struck{de \ill{}} plus d'un ton qu'avant cette
addition; car il étoit entre mi et fa 1.

A l'égard dela figure, j'ai vû la ligne nodale parallelle
au côté chargé se raprocher de ce côté, desorte qu'en
suivant cette indication, il falloit aussi remonter les
doigts, afin d'obtenir la figure aussi nette qu'elle pût
se former. La ligne nodale parallelle au côté chargé,
s'est alors formée à 45 millimetres de ce côté; elle
occupoit près de deux millimetres, parceque, comme
on sait, les lignes nodales sont d'autant plus \struck{foibles}
larges que les sons sont plus foibles. Il y a toujours
environ 1 millimetre d'engagé dans la cire; ainsi
il restoit 58 millimetres entre la ligne nodale
et l'extrémité opposée au côté chargé.
\note{En tête, au milieu, « (130) », non barré ; aucun nombre en haut à
droite : c'est le dos du feuillet de la vue 539. La première ligne
achève la phrase de la vue 539 (« en prenant pour premier octave celui
auquel appartient / le \emph{la} du diapason »). « la » est souligné,
comme le $x$ de « l'axe des $x$ ». Le « (1) » qui suit « sol » et le
« 1 » qui suit « fa » sont, semble-t-il, le numéro de l'octave ainsi
définie, non un appel de note : aucune note n'est au bas de la page.
Devant « plus d'un ton », « de » et un mot court sont barrés ensemble ;
le second ne se lit pas. Le mot lu « mi » porte à la fin une petite
tache. « foibles » est barré à la fin de la ligne, et récrit à la ligne
suivante à sa place. Le « 4 » de « 45 » est tracé comme un « h » ; sous
le chiffre, trois petites boucles sont des taches. Le texte s'arrête
sur un point aux deux tiers de la page ; le bas du feuillet est blanc.
La vue 541, hors de ce lot, commence un nouvel alinéa en tête de la page
(131).}

\end{document}
